
% ===== Slide 32 =====
\section{The Matrix Equation}
\begin{frame}{The Matrix Equation}
  \centering
  \Large 1.4 The Matrix Equation Ax=b

  \bigskip
  \Huge $A\mathbf{x}=\mathbf{b}$
\end{frame}

% ===== Slide 33 =====
\begin{frame}{Two Ways to Denote an $m\times n$ Matrix}
  In terms of the entries of $A$:
  \[
    A=\begin{bmatrix}
      a_{11}&\cdots&a_{1j}&\cdots&a_{1n}\\
      \vdots&&\vdots&&\vdots\\
      a_{i1}&\cdots&a_{ij}&\cdots&a_{in}\\
      \vdots&&\vdots&&\vdots\\
      a_{m1}&\cdots&a_{mj}&\cdots&a_{mn}
    \end{bmatrix}.
  \]
  In terms of its columns:
  \[
    A=\begin{bmatrix}\mathbf{a}_1&\mathbf{a}_2&\cdots&\mathbf{a}_n\end{bmatrix},
    \qquad \mathbf{a}_i\in\mathbb{R}^m\quad(i=1,\ldots,n).
  \]

  
\end{frame}

\begin{frame}{Two Ways to Denote a matrix}
   
 In terms of its columns:
  \[
    A=\begin{bmatrix}\mathbf{a}_1&\mathbf{a}_2&\cdots&\mathbf{a}_n\end{bmatrix},
    \qquad \mathbf{a}_i\in\mathbb{R}^m\quad(i=1,\ldots,n).
  \]

   For example:

   \[
     \mathbf{a}_1 = \begin{bmatrix}
       a_{11}\\\vdots\\a_{i1}\\\vdots\\a_{m1} 
     \end{bmatrix}.
  \]

\end{frame}

% ===== Slide 34 =====
\begin{frame}{Definition of the Matrix--Vector Product}
  If $A$ is an $m\times n$ matrix with columns
  $\mathbf{a}_1,\ldots,\mathbf{a}_n$, and
  $\mathbf{x}\in\mathbb{R}^n$, then
  \[
    A\mathbf{x}
    =\begin{bmatrix}\mathbf{a}_1&\mathbf{a}_2&\cdots&\mathbf{a}_n\end{bmatrix}
     \begin{bmatrix}x_1\\x_2\\\vdots\\x_n\end{bmatrix}
    =x_1\mathbf{a}_1+x_2\mathbf{a}_2+\cdots+x_n\mathbf{a}_n.
  \]
  Thus $A\mathbf{x}$ is the \textbf{linear combination} of the columns of $A$ using
  the corresponding entries of $\mathbf{x}$ as weights.

  \medskip
  $A\mathbf{x}$ is defined only when the number of columns of $A$ equals
  the number of entries of $\mathbf{x}$. ("Length" of $A$ = "Height" of x)
  
  Moreover,
  $\mathbf{a}_i\in\mathbb{R}^m$ and $A\mathbf{x}\in\mathbb{R}^m$.
\end{frame}

% ===== Slide 35 =====
\begin{frame}{Example: Computing $A\mathbf{x}$ by Columns}
  \[
    \underbrace{\begin{bmatrix}1&-4\\3&2\\0&5\end{bmatrix}}_{A}
    \underbrace{\begin{bmatrix}7\\-6\end{bmatrix}}_{\mathbf{x}}
    =7\begin{bmatrix}1\\3\\0\end{bmatrix}
      -6\begin{bmatrix}-4\\2\\5\end{bmatrix}
    =\begin{bmatrix}7\\21\\0\end{bmatrix}
      +\begin{bmatrix}24\\-12\\-30\end{bmatrix}
    =\underbrace{\begin{bmatrix}31\\9\\-30\end{bmatrix}}_{\mathbf{b}}.
  \]

  Here, $\mathbf{A}$
\end{frame}

% ===== Slide 36 =====
\begin{frame}{Example: Write a Linear Combination as $A\mathbf{x}$}
  For $\mathbf{v}_1,\mathbf{v}_2,\mathbf{v}_3\in\mathbb{R}^m$, write the linear combination 
  $3\mathbf{v}_1-5\mathbf{v}_2+7\mathbf{v}_3$ as a matrix times a vector.\pause

  \medskip
  Answer: place $\mathbf{v}_1,\mathbf{v}_2,\mathbf{v}_3$ into the columns of $A$
  and the weights $3,-5,7$ into $\mathbf{x}$:\pause
  \[
    3\mathbf{v}_1-5\mathbf{v}_2+7\mathbf{v}_3
    =\pause\begin{bmatrix}\mathbf{v}_1&\mathbf{v}_2&\mathbf{v}_3\end{bmatrix}
      \begin{bmatrix}3\\-5\\7\end{bmatrix}
    =A\mathbf{x}.
  \]
\end{frame}

% ===== Slide 37 =====
\begin{frame}{Three Forms of the Same Linear System}
  \[
    \begin{cases}
      2x_1+3x_2+4x_3=9,\\
      -3x_1+x_2=-2.
    \end{cases}
  \]
  Vector equation:
  \[
    x_1\begin{bmatrix}2\\-3\end{bmatrix}
    +x_2\begin{bmatrix}3\\1\end{bmatrix}
    +x_3\begin{bmatrix}4\\0\end{bmatrix}
    =\begin{bmatrix}9\\-2\end{bmatrix}.
  \]
  Matrix equation:
  \[
    \begin{bmatrix}2&3&4\\-3&1&0\end{bmatrix}
    \begin{bmatrix}x_1\\x_2\\x_3\end{bmatrix}
    =\begin{bmatrix}9\\-2\end{bmatrix},
    \qquad A\mathbf{x}=\mathbf{b}.
  \]
\end{frame}

% ===== Slide 38 =====
\begin{frame}{Three Equivalent Views of a Linear System}
  A linear system may be viewed
  \begin{enumerate}
    \item as a system of linear equations;
    \item as a vector equation
      \[
        x_1\mathbf{a}_1+x_2\mathbf{a}_2+\cdots+x_n\mathbf{a}_n=\mathbf{b};
      \]
    \item as a matrix equation $A\mathbf{x}=\mathbf{b}$.
  \end{enumerate}
\end{frame}

% ===== Slide 39 =====
\begin{frame}{Theorem 3}
  If $A$ is an $m\times n$ matrix with columns
  $\mathbf{a}_1,\ldots,\mathbf{a}_n$ and $\mathbf{b}\in\mathbb{R}^m$, then
  the matrix equation
  \[
    A\mathbf{x}=\mathbf{b}
  \]
  has the same solution set as the vector equation
  \[
    x_1\mathbf{a}_1+x_2\mathbf{a}_2+\cdots+x_n\mathbf{a}_n=\mathbf{b},
  \]
  which has the same solution set as the linear system with augmented matrix
  \[
    \begin{bmatrix}\mathbf{a}_1&\mathbf{a}_2&\cdots&\mathbf{a}_n&\mathbf{b}\end{bmatrix}.
  \]
  Consequently, $A\mathbf{x}=\mathbf{b}$ has a solution if and only if
  $\mathbf{b}$ is a linear combination of the columns of $A$.
\end{frame}

% ===== Slide 40 =====
\begin{frame}{Theorem 4 and Example}
  For an $m\times n$ matrix $A$, the following are logically equivalent:
  \begin{enumerate}
    \item For every $\mathbf{b}\in\mathbb{R}^m$, the equation
          $A\mathbf{x}=\mathbf{b}$ has a solution.
    \item Every $\mathbf{b}\in\mathbb{R}^m$ is a linear combination of the
          columns of $A$.
  \end{enumerate}
  
\end{frame}

\begin{frame}{How to determine an equation is consistent for all vectors}
  EXAMPLE: Let
  \[
    A=\begin{bmatrix}1&4&5\\-3&-11&-14\\2&8&10\end{bmatrix},\qquad
    \mathbf{b}=\begin{bmatrix}b_1\\b_2\\b_3\end{bmatrix},
  \]

  Is the equation $A\mathbf{x}=\mathbf{b}$ consistent for all $\mathbf{b}$?

  Answer: augmented matrix corresponding to $A\mathbf{x}=\mathbf{b}$:
  
  \[
    \left[\begin{array}{cccc}1&4&5&b_1\\-3&-11&-14&b_2\\2&8&10&b_3\end{array}\right]
    \sim
    \left[\begin{array}{cccc}
      1&4&5&b_1\\
      0&1&1&3b_1+b_2\\
      0&0&0&-2b_1+b_3
    \end{array}\right].
  \]

  (Can you find a value of $b_1$ and $b_3$ that makes $-2b_1+b_3$ non zero?)

  Hence $A\mathbf{x}=\mathbf{b}$ is \alert{not} consistent for every $\mathbf{b}$.
\end{frame}
% ===== Slide 41 =====
\begin{frame}{Exercise: Consistency for Every $\mathbf{b}$?}
  Compute $A\mathbf{x}$ where 
  \[
    A=\begin{bmatrix}1&2\\3&4\\5&6\end{bmatrix},\qquad
    \mathbf{b}=\begin{bmatrix}b_1\\b_2\\b_3\end{bmatrix}.
  \]
  Is $A\mathbf{x}=\mathbf{b}$ consistent for every possible $\mathbf{b}$?
  \[
    \left[\begin{array}{cc|c}1&2&b_1\\3&4&b_2\\5&6&b_3\end{array}\right]
    \sim
    \left[\begin{array}{cc|c}
      1&0&-2b_1+b_2\\
      0&1&\dfrac{3b_1-b_2}{2}\\
      0&0&b_1-2b_2+b_3
    \end{array}\right].
  \]
  The equation is not consistent for every $\mathbf{b}$ because some choices
  make $b_1-2b_2+b_3\ne0$.
\end{frame}

% ===== Slide 42 =====
\begin{frame}{Properties of the Matrix--Vector Product $A\mathbf{x}$}
  \begin{block}{Theorem 5}
    If $A$ is an $m\times n$ matrix,
    $\mathbf{u},\mathbf{v}\in\mathbb{R}^n$, and $c$ is a scalar, then
    \begin{enumerate}
      \item $A(\mathbf{u}+\mathbf{v})=A\mathbf{u}+A\mathbf{v}$;
      \item $A(c\mathbf{u})=cA\mathbf{u}$.
    \end{enumerate}
  \end{block}
\end{frame}

% ===== Slide 43 =====
\begin{frame}{Computation of $A\mathbf{x}$}
  Let
  \[
    A=\begin{pmatrix}2&3&4\\-1&5&-3\\6&-2&8\end{pmatrix},\qquad
    \mathbf{x}=\begin{pmatrix}x_1\\x_2\\x_3\end{pmatrix}.
  \]
  From the definition,
  \begin{align*}
    A\mathbf{x}
      &=x_1\begin{pmatrix}2\\-1\\6\end{pmatrix}
       +x_2\begin{pmatrix}3\\5\\-2\end{pmatrix}
       +x_3\begin{pmatrix}4\\-3\\8\end{pmatrix}\\
      &=\begin{pmatrix}
          2x_1+3x_2+4x_3\\
          -x_1+5x_2-3x_3\\
          6x_1-2x_2+8x_3
        \end{pmatrix}.
  \end{align*}
\end{frame}

% ===== Slide 44 =====
\begin{frame}{Computing $A\mathbf{x}$ Row by Row}
  \[
    \begin{pmatrix}2&3&4\\-1&5&-3\\6&-2&8\end{pmatrix}
    \begin{pmatrix}x_1\\x_2\\x_3\end{pmatrix}
    =\begin{pmatrix}
       \alert{2x_1+3x_2+4x_3}\\
       -x_1+5x_2-3x_3\\
       6x_1-2x_2+8x_3
     \end{pmatrix}.
  \]
  Each output entry is the dot product of one row of $A$ with $\mathbf{x}$:
  \[
    \begin{aligned}
      (2,3,4)\mathbf{x}&=2x_1+3x_2+4x_3,\\
      (-1,5,-3)\mathbf{x}&=-x_1+5x_2-3x_3,\\
      (6,-2,8)\mathbf{x}&=6x_1-2x_2+8x_3.
    \end{aligned}
  \]
\end{frame}

% ===== Slide 45 =====
\begin{frame}{Row--Vector Rule for Computing $A\mathbf{x}$}
  If $A\mathbf{x}$ is defined, then its $i$th entry is the sum of the products
  of corresponding entries from row $i$ of $A$ and from $\mathbf{x}$.
  \begin{align*}
    \begin{bmatrix}1&2&-1\\0&-5&3\end{bmatrix}
      \begin{bmatrix}4\\3\\7\end{bmatrix}
      &=\begin{bmatrix}3\\6\end{bmatrix},\\
    \begin{bmatrix}2&-3\\8&0\\-5&2\end{bmatrix}
      \begin{bmatrix}4\\7\end{bmatrix}
      &=\begin{bmatrix}-13\\32\\-6\end{bmatrix},\\
    \begin{bmatrix}1&0&0\\0&1&0\\0&0&1\end{bmatrix}
      \begin{bmatrix}r\\s\\t\end{bmatrix}
      &=\begin{bmatrix}r\\s\\t\end{bmatrix}.
  \end{align*}
\end{frame}

% ===== Slide 46 =====
\begin{frame}{Identity Matrix}
  % Chinese subtitle in source: 单位矩阵（单位阵）
  The $n\times n$ identity matrix is
  \[
    I=I_n=
    \begin{pmatrix}
      1&0&\cdots&0\\
      0&1&\cdots&0\\
      \vdots&\vdots&\ddots&\vdots\\
      0&0&\cdots&1
    \end{pmatrix}.
  \]
  Its diagonal entries are $1$ and all other entries are $0$. For example,
  \[
    I_3\begin{bmatrix}r\\s\\t\end{bmatrix}
    =\begin{bmatrix}r\\s\\t\end{bmatrix}.
  \]
  In general,
  \[
    I_n\mathbf{x}=\mathbf{x}\qquad\text{for every }\mathbf{x}\in\mathbb{R}^n.
  \]
\end{frame}
